% from ega4-1-fr.tex
\oldpage[IV-1]{246}

où $X'_\lambda=\operatorname{Spec}(A'_\lambda)$ est affine, et $f_\lambda$
est un morphisme $X'_\lambda\to X$~; alors, avec les notations de (1.9.3.2),
on a par hypothèse $f(X')=\varnothing$, donc $X'=\varnothing$~; mais cela
entraîne par (1.9.2, (i)) qu'il existe une partie finie $J$ de $L$ telle que
$X'_J=\varnothing$, donc
$f_J(X'_J)=\bigcap_{\lambda\in J}F_\lambda=\varnothing$.
\end{proof}

\begin{corollary}[1.9.9]
\label{IV.1.9.9-fr}
Soient $X$ un préschéma quasi-compact, $F$ une partie pro-constructible de
$X$, $(E_\lambda)_{\lambda\in L}$ une famille de parties ind-constructibles
de $X$ telle que $F\subset\bigcup_{\lambda\in L}E_\lambda$~; alors il existe
une partie finie $J$ de $L$ telle que
$F\subset\bigcup_{\lambda\in J}E_\lambda$. En particulier, de tout
recouvrement de $X$ par des parties ind-constructibles, on peut extraire un
recouvrement fini de $X$.
\end{corollary}

\begin{proof}
Cela résulte de (1.9.7) par passage aux complémentaires.
\end{proof}

\begin{proposition}[1.9.10]
\label{IV.1.9.10-fr}
Soit $X$ un préschéma. Pour qu'une partie $E$ de $X$ soit
ind-constructible, il est nécessaire que, pour tout $x\in X$, l'intersection
$E\cap\overline{\{x\}}$ soit un voisinage de $x$ dans
$\overline{\{x\}}$. Si $X$ est localement noethérien, cette condition est
suffisante.
\end{proposition}

\begin{proof}
Supposons $E$ ind-constructible, et soit $Y$ l'intersection de
$\overline{\{x\}}$ et d'un ouvert affine dans $X$, contenant $x$~; il y a
donc un sous-préschéma de $X$ ayant $Y$ pour espace sous-jacent (I, 5.2.1),
et si $j:Y\to X$ est l'injection canonique,
$E\cap Y=j^{-1}(E)$ est ind-constructible dans $Y$ (1.9.5, (vi)). Comme $Y$
est quasi-compact et séparé, $E\cap Y$ est donc réunion de parties
constructibles de $Y$ (1.9.5, (iii))~; par suite, il y a deux ouverts $U$,
$V$ dans $Y$ tels que
$x\in U\cap\complement V\subset E\cap Y$ (0\textsubscript{III}, 9.1.3).
Mais comme $x$ est point générique de l'espace irréductible $Y$, $V$ est
nécessairement vide et l'on a $U\subset E\cap Y$.

Supposons maintenant $X$ localement noethérien, et soit $E$ une partie de
$X$ vérifiant la condition de l'énoncé. En vertu de la définition (1.9.4),
on peut se borner au cas où $X$ est noethérien. Alors, pour tout $x\in E$,
$E\cap\overline{\{x\}}$ contient une partie non vide de la forme
$U\cap\overline{\{x\}}$, où $U$ est ouvert dans $X$~; cela montre que $E$
est réunion de parties localement fermées de $X$, donc est
ind-constructible (1.9.6).
\end{proof}

\begin{proposition}[1.9.11]
\label{IV.1.9.11-fr}
Soient $X$ un préschéma, $E$ une partie de $X$. Les propriétés suivantes sont
équivalentes~:
\begin{enumerate}
  \item[a)] $E$ est localement constructible.
  \item[b)] $E$ est ind-constructible et pro-constructible.
  \item[c)] $E$ et $X-E$ sont pro-constructibles.
  \item[d)] $E$ et $X-E$ sont ind-constructibles.
\end{enumerate}
\end{proposition}

\begin{proof}
Il suffit évidemment de prouver que $d)$ entraîne $a)$, et on peut se borner
au cas où $X$ est affine. Alors (1.9.5, (iii)), $E$ (resp. $X-E$) est réunion
d'une famille $(E_\lambda)$ (resp. $(E'_\mu)$) de parties constructibles de
$X$~; comme les $E_\lambda$ et les $E'_\mu$ forment un recouvrement de $X$,
il résulte de (1.9.9) qu'il y a des indices en nombre fini $\lambda_i$,
$\mu_j$ tels que les $E_{\lambda_i}$ et les $E'_{\mu_j}$ forment un
recouvrement de $X$~; cela implique que $E$ est réunion des
$E_{\lambda_i}$, donc est constructible.
\end{proof}

\begin{corollary}[1.9.12]
\label{IV.1.9.12-fr}
Soit $f:X\to Y$ un morphisme surjectif, qui est, soit quasi-compact, soit
localement de présentation finie. Pour qu'une partie $E$ de $Y$ soit
localement constructible (resp. pro-constructible, resp. ind-constructible),
il faut et il suffit que $f^{-1}(E)$ le soit dans $X$.
\end{corollary}

